%! Functions — Unit 1, Lesson 1.3 — Lesson Plan
\documentclass[10pt]{article}
\usepackage{algebra2-article}
\usepackage{algebra2-boxes}
\usepackage{graphicx}   % -article does NOT load graphicx; needed for thumbnails/screenshots
\graphicspath{{images/}}
\newcommand{\TallMath}[1]{$\displaystyle #1\rule[-1.4em]{0pt}{3.2em}$}
\newcommand{\CourseName}{Algebra 2: Shepherd}
\newcommand{\SchoolYear}{2026--2027}
\newcommand{\MeetingLength}{55 minutes}
\newcommand{\UnitNumberName}{Unit 1: Foundations --- Algebra 1 Review \quad}
\newcommand{\LessonNumberName}{Lesson 1.3: Functions}

\begin{document}
\begin{center}
    {\Large\bfseries \CourseName: \SchoolYear \\
    {\small\bfseries \UnitNumberName \LessonNumberName}}
\end{center}

\begin{tcolorbox}[colback=forestbg, colframe=forest, boxrule=0.9pt, arc=2mm,
    left=2mm, right=2mm, top=1.5mm, bottom=1.5mm]
    \textbf{Primary Objective:} Students can decide whether a relation --- given as ordered pairs, a
    table, a mapping, or a graph --- is a \emph{function}, justify the decision with the vertical
    line test, and state its domain and range; use function notation in \emph{both} directions,
    finding $f(x)$ from $x$ and $x$ from $f(x)$, algebraically and graphically; determine a linear
    function's domain, range, \emph{zeros}, slope, and intercepts and interpret each in context;
    write the equation of a line in slope-intercept, point-slope, and standard form --- from a graph,
    from two points, from a slope and a point, vertically, horizontally, and parallel or
    perpendicular to a given line; explain how $k\,f(x)$ and $f(x)+k$ transform the parent $y=x$;
    read the key characteristics of quadratic and exponential functions off their graphs; and
    compare and contrast the three families using tables and graphs.
    \\[2pt]
    \textbf{Standards (2023 VA SOL --- Algebra 1, reviewed):} this lesson reviews the
    \textbf{Functions} domain, the third and last domain area of Unit 1. \textbf{A.F.1} --- domain,
    range, zeros, slope, and intercepts of a linear function presented algebraically or graphically,
    including interpretation in context (a); how transformations to the parent $y=x$ affect slope
    and $y$-intercept (b); equivalent algebraic forms --- slope-intercept, standard, point-slope ---
    and what each reveals (c); writing the equation of a line from a graph, from two points, from a
    slope and a point, and as $x=a$ or $y=c$ (d); parallel and perpendicular lines through a given
    point (e); graphing a linear function, including contextual ones (f); determining $f(x)$ from
    $x$ and $x$ from $f(x)$ (g); and comparing linear functions across algebraic, graphical,
    tabular, and contextual representations (h). \textbf{A.F.2} --- deciding whether a relation is a
    function and giving its domain and range (a); key characteristics of a quadratic --- zeros,
    $y$-intercept, vertex as maximum or minimum, domain and range including when restricted by
    context (b); graphing a quadratic, including by $f(x)+k$ and $kf(x)$ (c); connecting standard
    and factored forms to the graph (d); key characteristics of an exponential $y=ab^x$ (e);
    graphing an exponential (f); evaluating and solving for quadratic and exponential functions and
    explaining the meaning of $x$ and $f(x)$ in context (g); and comparing and contrasting the three
    families using tables and graphs (h).
    \\[2pt]
    \textbf{Unit 1's ``characteristics'' row lands here.} Unit 1 is the only unit whose Lesson 0 is
    not a characteristics lesson, so this lesson carries the course spine's opening row --- function
    definition and notation, and domain and range read from a graph. Everything Unit 2 does to a
    line, Unit 3 to a parabola, and Unit 7 to a growth curve starts from the vocabulary built today.
    \\[2pt]
    \textbf{Prerequisite:} Lessons 1.0--1.2. From 1.2 especially: solving $f(x)=0$ is how a zero
    gets found, so the whole ``backward'' half of function notation is last lesson's work under a
    new name. From 1.2's literal equations: the three \emph{forms} of a line are one equation solved
    for different things. From 1.1: factoring is what makes a quadratic's zeros visible.
\end{tcolorbox}

\renewcommand{\arraystretch}{1.4}
\begin{skillbox}[Priority Ideas \& Skills]{goldbox}
\begin{tabularx}{\linewidth}{@{}X|X@{}}
\textbf{Priority Ideas \& Skills} & \textbf{Key Understandings (the ``why'')} \\[2pt]
\hline
\vspace{-6pt}
\begin{itemize}[leftmargin=1.1em, nosep, after=\vspace{2pt}]
  \item \textbf{Decide} whether a relation is a function from pairs, a table, a mapping, or a graph; apply the \textbf{vertical line test}.
  \item \textbf{State} the domain and range of a relation, including when \emph{context} restricts them.
  \item \textbf{Evaluate} $f(x)$ for a given $x$, and \textbf{solve} for $x$ given $f(x)$ --- from a rule and from a graph.
  \item \textbf{Find} a linear function's zeros, intercepts, and slope, and \textbf{interpret} each in context with units.
  \item \textbf{Write} a line's equation in all three forms, from a graph, two points, a slope and a point, and as $x=a$ or $y=c$.
  \item \textbf{Write} equations of parallel and perpendicular lines.
  \item \textbf{Explain} how $k\,f(x)$ and $f(x)+k$ change the parent $y=x$.
  \item \textbf{Read} vertex, zeros, intercepts, domain, and range off a parabola; \textbf{read} starting value, growth or decay, and range off an exponential.
  \item \textbf{Identify} a family from a bare table using differences and ratios.
\end{itemize}
&
\vspace{-6pt}
\begin{itemize}[leftmargin=1.1em, nosep, after=\vspace{2pt}]
  \item An equation was a question about \emph{numbers}. A function is a \textbf{machine}, and its graph shows every answer at once.
  \item The promise --- one input, exactly one output --- constrains \textbf{inputs only}. Repeated outputs are fine; that is why $y=-7$ is a function and $x=-7$ is not.
  \item The vertical line test is not a trick. It is the definition, \emph{drawn}.
  \item \emph{Zero}, \emph{solution of $f(x)=0$}, and \emph{$x$-intercept} are three names for one thing --- which is why Lesson 1.2 was really about functions all along.
  \item The three forms are not three facts. They are one equation solved for different things, and each is honest about a different feature.
  \item Slope is a \textbf{rate}, and in context it always has units: gallons per minute, dollars per shirt.
  \item Context restricts a domain. The algebra permits every real number; the situation does not, and the situation wins.
  \item Linear adds, quadratic turns around, exponential \emph{multiplies} --- which is why the exponential has no zero and why it eventually wins any race.
\end{itemize}
\end{tabularx}
\end{skillbox}

\begin{skillbox}[{Vocabulary, Concepts, \& Theorems}]{sky}
\renewcommand{\arraystretch}{1.3}
\begin{tabularx}{\linewidth}{@{}l X@{}}
\textbf{Relation} & Any pairing of inputs with outputs --- ordered pairs, a table, a mapping, or a graph. \\
\textbf{Function} & A relation in which every input has \emph{exactly one} output. Two inputs may share an output; one input may not have two. \\
\textbf{Vertical line test} & A graph is a function exactly when no vertical line crosses it more than once. It is the definition drawn: a vertical line is one input. \\
\textbf{Domain / range} & The set of all inputs / the set of all outputs. Either may be restricted by context. \\
\textbf{Function notation} & $f(x)$, read ``$f$ of $x$'': the output of $f$ at input $x$. Not multiplication. ``$f(3)=0$'' $=$ ``$(3,0)$ is on the graph.'' \\
\textbf{Zero of a function} & An input making the output $0$ --- a solution of $f(x)=0$, and an $x$-intercept of the graph. \\
\textbf{Slope / rate of change} & \TallMath{m=\frac{y_2-y_1}{x_2-x_1}} --- the change in output per unit change in input. Constant exactly for linear functions. \\
\textbf{Slope-intercept form} & $y=mx+b$. Shows the slope and the $y$-intercept immediately. \\
\textbf{Point-slope form} & $y-y_1=m(x-x_1)$. Shows a slope and one point on the line. \\
\textbf{Standard form} & $Ax+By=C$. Gives both intercepts quickly; the only form that can also write $x=a$. \\
\textbf{Parallel / perpendicular} & Equal slopes / opposite reciprocal slopes (their product is $-1$). \\
\textbf{Parent function} & The simplest member of a family: $y=x$, $y=x^2$, $y=b^x$. \\
\textbf{Transformation} & $k\,f(x)$ stretches or reflects (for a line, it changes the slope); $f(x)+k$ shifts vertically (it changes the $y$-intercept). \\
\textbf{Vertex} & The turning point of a parabola --- the maximum or minimum. Its input answers ``where,'' its output answers ``how much.'' \\
\textbf{Exponential function} & $y=ab^x$: $a$ is the starting value ($b^0=1$), $b$ is the constant multiplier. $b>1$ grows, $0<b<1$ decays. \\
\textbf{Constant difference / ratio} & Linear data have a constant first difference; quadratic, a constant \emph{second} difference; exponential, a constant ratio. \\
\end{tabularx}
\end{skillbox}

\begin{fixedskillbox}[Activate Prior Knowledge \& Spiral Review]{forestbg}
    The Warm-Up is the lesson in miniature: \textbf{one rule, $2x-6$, asked three ways.}
    \textbf{(1)} Solve $2x-6=0$ and $x^2-2x-3=0$ --- pure Lesson 1.2, and the answers are about to
    be renamed \emph{zeros}. \textbf{(2)} Evaluate $2x-6$, $x^2-2x-3$, and $2^x$ at $x=0$ and
    $x=3$ --- pure Lesson 1.1 --- with two planted questions: which rules gave zero at $x=3$ (the
    first two, matching item 1 exactly), and whether \emph{any} input makes $2^x$ zero (no, and
    students should try negatives to convince themselves). \textbf{(3)} Read a pre-drawn line for
    its intercepts, direction, and step size --- which is slope arrived at by counting boxes.
    The closing question asks how students can tell that item 1(a), item 2's first row, and the
    graph are the \emph{same rule}. Take that answer publicly: the representations agree because
    they are one object seen three ways, which is \textbf{A.F.1h} stated by the class rather than by
    the teacher. Debrief the $2^x$ row without naming the asymptote --- that word belongs to Unit 5.
\end{fixedskillbox}

\begin{skillbox}[Hook]{forestbg}
    \textbf{Two beats, five minutes.} First, put three relations on the board and ask only whether
    each one gives every input exactly one output: \emph{(a)} $\{(1,3),(2,5),(4,9)\}$;
    \emph{(b)} $\{(1,3),(2,5),(1,7)\}$; \emph{(c)} $\{(1,6),(2,6),(4,6)\}$. Everyone rejects (b)
    correctly. The item that does the work is \emph{(c)} --- expect several students to reject it
    too, because ``the $6$ repeats.'' Do not correct it yet; collect a show of hands and write both
    counts on the board. Then state the promise once, precisely: \textbf{every input gets exactly
    one output}, and note what it does \emph{not} say --- nothing forbids two inputs from sharing an
    output. (c) is a function. That misconception, left alive, makes every parabola a non-function
    in Unit 3, so it is worth the ninety seconds now.
    \\[3pt]
    Second beat, and the reason this lesson has three families in it: put $f(x)=x$, $g(x)=x^2$, and
    $h(x)=2^x$ in one table for $x=0,1,2,3,4,5$ and have the room fill it. Then ask who is biggest
    at $x=2$ --- $g$ and $h$ tie at $4$. At $x=4$ --- they tie \emph{again}, at $16$. At $x=5$ ---
    $32$ beats $25$, and the exponential is gone for good. \textbf{The answer changed, twice.} Land
    the frame: \emph{today we ask one fixed set of questions --- domain, range, zeros, intercepts,
    rate --- of three different families, and the interesting part is how differently each one
    answers.} Leave the completed table on the board all period; Section 5 returns to it.
\end{skillbox}

\begin{skillbox}[Lesson]{forestbg}
\begin{multicols}{2}
\textbf{1. What makes a rule a function.} Open by naming the contrast with 1.2: an equation asks
about numbers, a function is a \emph{machine}. Give the definition, then show it surviving four
different costumes --- ordered pairs, table, mapping, graph --- and make the point that the test is
always the same test, just wearing different clothes (\textbf{A.F.2a}). Derive the \textbf{vertical
line test} rather than announcing it: a vertical line is the set of points with one fixed $x$, so
two hits means one input with two outputs. Use the three pre-drawn panels; the middle one is $x=2$,
a perfectly good \emph{line} that is not a function, and it pairs with $y=-2$ later in the period.
Then domain and range (a list for a set of pairs, ``all real numbers'' for an unbroken line), and
\textbf{function notation}: $f$ is the name, $x$ the input, $f(x)$ the output --- and ``$f(3)=0$''
is the same sentence as ``$(3,0)$ is on the graph.'' Close the section with the forward/backward
table (\textbf{A.F.1g}) and the line that ties Unit 1 together: \emph{forward is arithmetic;
backward is Lesson 1.2}. The $f(x)=0$ box has a name --- \textbf{zero} --- and it is the
$x$-intercept.

\textbf{2. Reading a line.} Take $f(x)=2x-6$, the Warm-Up's rule, and interrogate it: domain, range,
zero, $x$-intercept, $y$-intercept, slope (\textbf{A.F.1a}). Ask which two rows are the same fact
twice; that is the point of the table. Then slope as a \emph{rate}, computed from two points read
off the graph, and said in words --- ``up $2$ every time $x$ goes up $1$.'' Finish in context with
the draining pool $V(t)=800-25t$: the intercept is the starting volume, the slope is
$-25$ gallons per minute, the zero is when it is empty, and the domain is $0\le t\le32$ because
negative time means nothing and after $32$ minutes the model has stopped describing anything. Say
that explicitly --- the algebra allows every real number; the situation does not, and the situation
wins.

\columnbreak

\textbf{3. Writing the line.} Three forms, each honest about something different
(\textbf{A.F.1c}); rewrite $y-4=2(x-5)$ into the other two and ask what changed --- nothing, and the
skill used was 1.2's literal equations. Then four writing situations under one strategy, \emph{get a
slope and get a point}: slope-and-point, two points, vertical $x=a$, horizontal $y=c$
(\textbf{A.F.1d}). Ask which of the last two is a function and why; that closes the loop with
Section 1. Parallel and perpendicular next (\textbf{A.F.1e}), then the parent $y=x$ and its two
transformations (\textbf{A.F.1b}): $k\,f(x)$ moves the slope, $f(x)+k$ moves the $y$-intercept, and
each leaves the other alone. Match the three pre-drawn panels. Say out loud that this pair of moves
returns in Unit 2 on a V, Unit 3 on a parabola, and Unit 7 on a growth curve --- it is the single
most reused idea in the course.

\textbf{4. The other two families.} Same six questions, different answers. For
$g(x)=x^2-2x-3$ --- the curve students already met in 1.2 --- read zeros, $y$-intercept, vertex,
max/min, domain, and range off the graph, then connect $(x-3)(x+1)$ to the two intercepts
(\textbf{A.F.2b/d}). The new idea is that the \emph{range stops}, because the graph turns around.
For $h(x)=2^x$: $y$-intercept $(0,1)$, no $x$-intercept at all, range $y>0$, always increasing
(\textbf{A.F.2e}). Get the ``never reaches zero'' argument from the Hook's table and from negative
inputs. Finish with the doubling colony $P(t)=50\cdot2^t$, where $a$ is the starting amount and $b$
is the multiplier (\textbf{A.F.2g}).

\textbf{5. Side by side.} Return to the Hook's table and look at \emph{how} each column grows: add
the same amount, add a growing amount whose growth is constant, or multiply by the same amount
(\textbf{A.F.2h}). Name the three signatures --- constant difference, constant second difference,
constant ratio --- and immediately use them to identify three bare tables. Close on the overtake: the
exponential tied twice and then left, so ``which is bigger'' is not a fair question without saying
\emph{where}.
\end{multicols}
\end{skillbox}

\begin{skillbox}[Active Monitoring]{redbox}
Circulate during the notes practice and Tier work. Watch for and correct:
\begin{itemize}[leftmargin=1.2em, nosep]
  \item calling a relation a non-function because an \emph{output} repeats --- the single most damaging error of the day; it kills every parabola in Unit 3;
  \item confusing $x=a$ with $y=c$: writing ``$y=5$'' for a vertical line through $(5,-8)$;
  \item reading $f(x)$ as $f$ times $x$, or writing $f(x+2)$ when asked for $f(x)+2$;
  \item stalling on the \emph{backward} column --- ``find $x$ when $f(x)=9$'' --- which is an equation-solving gap, not a functions gap; send those students to 1.2's notes, not today's;
  \item giving the zero as a $y$-value, or giving the $y$-intercept when asked for the $x$-intercept;
  \item computing slope as $\frac{\text{run}}{\text{rise}}$, or losing the sign on a falling line;
  \item slope from two points with the subtraction order flipped in only one of the two differences;
  \item taking the reciprocal but not the opposite sign for a perpendicular slope;
  \item stating a context domain as ``all real numbers'' when negative time or negative width is meaningless --- and, the reverse, restricting a domain that has no reason to be restricted;
  \item saying an exponential ``has an $x$-intercept because it keeps dropping'' --- approaching zero is not reaching it;
  \item checking only first differences and calling a quadratic table ``not linear, so exponential.''
\end{itemize}
Cold-call prompts: ``Is that a repeated input or a repeated output? Which one is illegal?'' ``Give me
that zero as a \emph{point}.'' ``What are the units of your slope?'' ``Same question, other
direction: what $x$ makes $f(x)=12$?'' ``Would a vertical line ever hit this twice?'' ``Differences
or ratios --- which did you check, and what did the other one say?''
\end{skillbox}

\begin{skillbox}[Group Work \& Differentiation]{redbox}
\textbf{One In, One Out --- Three Families, One Set of Questions} activity (tiered; mirrors the notes).
\begin{multicols}{3}
\textbf{Tier R --- Remediate.} Four blocks, one per notes section. Function-or-not with domain and
range, including a pre-drawn sideways parabola to fail the vertical line test on. A forward/backward
notation table for $f(x)=3x-12$, with the zero named and rewritten as a point. A line to interrogate
for intercepts, zero, slope, and equation. Then family identification from three bare tables. Send
here any student who called a repeated output illegal.

\columnbreak
\textbf{Tier A --- Approaching.} Mixed and unlabelled. Four equation-writing items (two points, a
slope and a point in two forms, parallel, perpendicular) closing with a deliberate trap --- which of
them is \emph{not} a function, answer \emph{none}, because every line with a slope passes the
vertical line test. Then a quadratic doing a job: a three-sided garden bed $A(x)=x(6-x)$, read for
zeros, vertex, max, and a \emph{restricted} domain, with both zeros interpreted as bad designs
rather than wrong answers. Then an exponential doing a job: a rumour tripling hourly, read for $a$,
$b$, a backward lookup, and whether the model ever reaches zero.

\columnbreak
\textbf{Tier E --- Extension.} Structure work. Build a function from conditions only --- linear,
$f(2)=7$, $f(5)=16$ --- then argue why the word \emph{linear} was load-bearing by producing a second
function through the same two points. A four-item ``spot the error'' set, one per misconception from
Active Monitoring, where the error must be \textbf{named} and sorted into definition errors and
procedure errors. The overtake table extended to $x=6$, with the two ties located and an argument
for why doubling must eventually win. Closes with the sharpest question on the page: what does a
\emph{horizontal} line test measure, given that $y=x^2$ fails it and is still a function, while
$x=3$ passes it and is not? \emph{(Do not name ``one-to-one'' --- that is Unit 6.)}
\end{multicols}
\end{skillbox}

\begin{skillbox}[Individual Work \& Assessment]{redbox}
\textbf{Exit Ticket} (independent, no notes): three items. \textbf{(1)} Function or not, for a set of
ordered pairs with a repeated \emph{input} and a table with repeated \emph{outputs}, plus the domain
and range of the one that qualifies and the rule stated in the student's own words --- the diagnostic
is whether ``repeated outputs'' is treated as a violation. \textbf{(2)} For $f(x)=-3x+12$: evaluate,
give the $y$-intercept and slope, then \emph{find the zero} by writing and solving $-3x+12=0$, and
say where that zero appears on the graph. A student who produces $x=4$ but cannot answer the second
half has the procedure without the picture. \textbf{(3)} An \textbf{SOL-style multiple-choice} item:
which equation represents a pre-drawn line (answer $y=-2x+4$), with distractors that each break
exactly one thing --- the sign of the slope, the sign of the intercept, or rise and run swapped.
Collect and sort into ``got it / definition / zero / slope.'' This is the last exit ticket before
the unit test, so the piles should set the review day's emphasis rather than tomorrow's reteach.
\end{skillbox}

\begin{skillbox}[Reinforcement \& Extension]{goldbox}
\textbf{Homework:} the whole domain area. Four function-or-not items including the $y=-7$ / $x=-7$
pair; a forward/backward notation table with the zero written as a point; a pre-drawn line read for
seven characteristics including a graphical $f(6)$ and a graphical backward lookup; one line written
in all three forms with a note on what each reveals; six equation-writing items (slope and point,
two points, horizontal, vertical, parallel, perpendicular) ending in ``which one is not a
function''; three parent-transformation panels to name; a pre-drawn downward parabola read for six
characteristics and confirmed by factoring; a depreciating-car exponential read for $a$, $b$, the
percentage kept and lost, and whether the value ever reaches zero; four bare tables to classify by
differences and ratios; and one written item resolving the repeated-input/repeated-output confusion
directly. \textbf{Extension:} determined vs.\ undetermined --- build a linear function from two
conditions, then show the rule collapses without the word ``linear''; write two different quadratics
with the same zeros and explain why they share a vertex $x$-coordinate; and race $L(x)=100x$ against
$E(x)=2^x$ to find where the exponential passes a hundred-fold head start, then diagnose the
classmate who says ``$100$ is bigger than $2$.''
\textbf{Preview:} \emph{the Unit 1 test}, and then \textbf{Unit 2 --- Linear Functions}, opening
with Lesson 2.0, \emph{Characteristics of Linear Functions}. Say the structure out loud: from here
on every unit opens by meeting a family and learning to read its graph before manipulating it, and
today's six questions are the permanent question list --- one that only grows.
\end{skillbox}
% ── Teacher notes (migrated from the component keys) ──
% One per component, titled for it. They live here rather than in the _key
% files so that every component is the same length blank and keyed.

\begin{teachernote}[Warm-Up]
\small
\textbf{This warm-up is the lesson in miniature.} The same rule $2x-6$ appears three times --- solved
algebraically, evaluated numerically, and read graphically --- and the answers agree every time. That
agreement \emph{is} \textbf{A.F.1h}: the representations are not three topics, they are one object
seen from three sides. Take the closing question out loud before starting the notes.

\textbf{Item 1 is Lesson 1.2 being cashed in, and today it gets a name.} The values students just
found are about to be called the \textbf{zeros} of a function. Say the definition once here so it
lands as a renaming rather than a new procedure: a zero is a solution of $f(x)=0$, which is the
same thing as an $x$-intercept of the graph.

\textbf{The $2^x$ row is a deliberate plant.} Students will try $x=0$, then negatives, and find
$2^{-1}=\frac12$, $2^{-3}=\frac18$ --- shrinking but never arriving. Do \emph{not} name the
asymptote (that debuts in Unit 5); just get the observation on the board. Section 4 of the notes
needs it, and it is the cleanest contrast between the exponential family and the other two, both of
which hit zero somewhere.

\textbf{Item 3 also pre-loads slope.} ``Every time $x$ increases by $1$, $y$ changes by $+2$'' is
the rate of change, arrived at by counting boxes rather than by a formula --- which is where the
notes pick it up.
\end{teachernote}

\begin{teachernote}[Guided Notes]
\small
\textbf{The Hook's two beats do two different jobs.} Part 1 is the definition and takes three
minutes. Part 2 is the whole reason this lesson has three families in it: the ``who is biggest''
answer \emph{changes}, twice, and students who have only ever seen $x^2$ grow fast are surprised.
Leave the completed table on the board all period --- Section 5 comes straight back to it.

\textbf{Section 1: the misconception to kill on sight} is ``repeated $y$-values mean it isn't a
function.'' Hook item (c) and the line $y=-2$ in Section 3 both exist to make the point. A student
who carries this error into Unit 3 will call every parabola a non-function.

\textbf{Section 1 is also where Unit 1 closes its loop.} Say the sentence out loud: \emph{Lesson 1.2
was never really about equations.} Solving $2x-6=0$ was finding a zero; solving $x^2-2x-3=0$ was
finding two. The backward column of the notation table is exactly A.EI.1b wearing new clothes, and
naming that is what makes the review feel cumulative instead of repetitive.

\textbf{No sketching anywhere in this lesson --- deliberately.} \textbf{A.F.1f} and
\textbf{A.F.2c/f} ask students to graph, and here that is done by \emph{matching} an equation to a
pre-drawn graph, completing a table of values, and reading features off a figure. If you want a
board demonstration of graphing by transformation, do it yourself with the three panels in Section 3
and let students name the move.

\textbf{Section 3's transformation table is a deposit, not a withdrawal.} It looks like a small
Algebra 1 item and it is the single most reused idea in the course: $k\,f(x)$ and $f(x)+k$ return in
Unit 2 (the V), Unit 3 (the parabola), and Unit 7 (growth curves), and \textbf{A2.F.1} is that
lens applied family by family. Name the units out loud.

\textbf{Section 4: do not name the asymptote.} The exponential's ``never reaches zero'' is an
observation students should own by the end of today, but the word \emph{asymptote} and the
horizontal-line-it-approaches formalism belong to Unit 5, where the characteristics spine introduces
them. Saying it early costs the Unit 5 lesson its punch and gains nothing here.

\textbf{Section 5 is A.F.2h, the clause most often skipped.} Constant difference / constant second
difference / constant ratio is the only tool that identifies a family from a bare table, and it is
routinely assessed. If time is short, cut the parallel/perpendicular pair in Section 3 --- not this.
\end{teachernote}

\begin{teachernote}[Group Activity]
\small
\textbf{Tier R Part 2 is the highest-leverage block on the page.} The forward column is arithmetic
and the backward column is Lesson 1.2, and putting them in one table is what makes ``a zero is a
solution of $f(x)=0$'' feel obvious rather than memorised. Anyone who can do the left column but
freezes on the right does not have a function problem --- they have an equation-solving problem, and
should be sent back to 1.2's notes rather than reread today's definition.

\textbf{Tier A Part 1 item 5 is a deliberate trap and worth debriefing publicly.} The answer is
\emph{none}: every line written as $y=mx+b$ passes the vertical line test. Only $x=a$ fails, and
$x=a$ can never be written in slope-intercept form because it has no slope. Students who answer
``item 3, because it's parallel to a steep line'' are guessing.

\textbf{Tier A Part 2's two zeros are the interpretation beat.} $x=0$ and $x=6$ are correct
solutions of $A(x)=0$ and useless as bed designs. Students trained to ``throw out the negative
answer'' need the sharper rule: you discard a solution when the \emph{context} forbids it, and you
must be able to say what the context objected to.

\textbf{Tier E Part 4 is the best thing on the page.} It forces students to state what the vertical
line test actually measures, and (b) is the item that separates memorised procedure from
understanding. Do not name ``one-to-one'' or ``inverse'' --- those belong to Unit 6 --- but if a
group gets there on their own, tell them they have found next year's question a year early.

\textbf{Tier E Part 1(d) is deliberately open.} Any curved function through $(2,7)$ and $(5,16)$
proves the point; accept a described curve, a sketchless argument, or a piecewise rule. The idea
being paid for is that \emph{two points determine a line} --- and only a line.
\end{teachernote}

\begin{teachernote}[Exit Ticket]
\footnotesize
\textbf{Answer 3: B.} The line falls $2$ for every $1$ right, so $m=-2$, and it crosses at $(0,4)$.
Each distractor breaks exactly one thing: \textbf{A} loses the sign of the slope, \textbf{C} the sign
of the intercept, \textbf{D} swaps rise and run. A student choosing D can read a graph but has the
ratio upside down --- a two-minute fix; a student choosing A is not attending to direction at all.

\textbf{Sort into four piles:} \emph{got it} / \emph{definition} (item 1 --- called (a) a function,
or called (b) \emph{not} one because outputs repeat) / \emph{zero} (item 2d/2e) / \emph{slope}
(item 3). The \emph{definition} pile costs the most later: a student who thinks repeated
\emph{outputs} disqualify a relation will call every parabola a non-function in Unit 3, and one
sentence fixes it --- the promise is about inputs only. The \emph{zero} pile is the Unit 1
through-line: 2(d) is Lesson 1.2's equation solving renamed, and 2(e) is the whole point of the
renaming --- the number you solved for is a \emph{place on the graph}.

\textbf{Unit 1 ends here,} so let the piles set the review day's emphasis rather than tomorrow's
reteach.
\end{teachernote}

\begin{teachernote}[Homework]
\small
\textbf{Grade problem 1 first, then problem 10.} They are the same question asked twice, and
together they diagnose the one misconception that does real damage: students who think a repeated
\emph{output} disqualifies a relation. Anyone who marked (a) or (c) ``not a function'' will call
every parabola a non-function in Unit 3 and every horizontal-line model a non-function in Unit 2.
One sentence fixes it --- the promise constrains inputs only --- and it is cheaper to say now.

\textbf{Problem 2 is the Unit 1 through-line, graded.} The left column is arithmetic; the right
column is Lesson 1.2. If a student can do the left and not the right, the gap is equation solving,
not function notation, and the Unit 1 test will expose it in three different places.

\textbf{Problem 5(d) and problem 1(d) are the same fact} arriving from opposite directions --- write
a vertical line, then be told it is not a function. If a student answers 5(d) as ``$y=5$'' they have
the vertical/horizontal pair backwards, which is worth thirty seconds of correction and is a
perennial SOL item.

\textbf{Problem 7's factoring check is where A.F.2d gets earned.} Reading zeros off a picture is
easy; producing $-(x-3)(x+1)$ and watching it agree is the connection the standard actually asks
for. Watch for the dropped leading negative --- students who factor $x^2-2x-3$ and forget the sign
get the right zeros for the wrong function and will not notice, because the zeros are unaffected.

\textbf{Problem 8 is the honest-modelling beat.} ``The car is never worth \$0'' is mathematically
correct and practically false, and saying so is not a digression --- \textbf{A.F.2e} asks students
to interpret characteristics \emph{in context}, which includes noticing where the model stops
describing reality. Accept any answer that separates what the function says from what a used car
lot would say.

\textbf{The Extension's item (3) is the best question on the page} and the one to put on the board
if anyone gets it. It also plants Unit 7: exponential growth is not ``fast growth,'' it is growth
whose \emph{rate} is proportional to what is already there, which is why no linear head start
survives it.

\textbf{Unit 1 ends here.} Use this set to build the review day: the four piles from the exit ticket
plus problems 1, 2, and 7 cover most of what the unit test will ask.
\end{teachernote}

\end{document}
